  % !TeX encoding = UTF-8
  % !TeX program = pdflatex
  % Arquivo independente: não requer imagens externas nem arquivo .bib.
  % Compilar duas vezes com pdfLaTeX para atualizar sumário e referências.
  % Antes da entrega: revisar funções, preencher URL do vídeo e commit final.
  \documentclass[12pt,a4paper]{article}
  \usepackage[T1]{fontenc}
  \usepackage[utf8]{inputenc}
  \usepackage[brazil]{babel}
  \usepackage{mathptmx}
  \usepackage{amsmath}
  \usepackage{geometry}
  \geometry{top=3cm,left=3cm,right=2cm,bottom=2cm}
  \usepackage{setspace}
  \onehalfspacing
  \usepackage{microtype}
  \usepackage{booktabs,tabularx,array}
  \usepackage{graphicx,float,caption,xcolor}
  \usepackage{enumitem,fancyhdr,titlesec}
  \usepackage{pgfplots}
  \pgfplotsset{compat=1.18}
  \usepackage{hyperref}
  \usepackage{url}
  \hypersetup{colorlinks=true,linkcolor=black,citecolor=black,urlcolor=blue,
   pdftitle={GenResearch-Servidor: memória, armazenamento, filas, cache e sincronização},
   pdfauthor={Equipe GenResearch - UFS}}
  \urlstyle{same}
  \setlength{\parindent}{1.5cm}
  \setlength{\parskip}{0pt}
  \setlength{\emergencystretch}{2em}
  \setlist{nosep,leftmargin=1.5cm}
  \captionsetup{font=small,labelfont=bf,justification=justified,singlelinecheck=false}
  \titleformat{\section}[block]{\normalfont\bfseries}{\thesection}{1em}{\MakeUppercase}
  \titleformat{\subsection}[block]{\normalfont\mdseries}{\thesubsection}{1em}{\MakeUppercase}
  \titlespacing*{\section}{0pt}{0.6\baselineskip}{0.6\baselineskip}
  \titlespacing*{\subsection}{0pt}{0.5\baselineskip}{0.3\baselineskip}
  \pagestyle{fancy}
  \setlength{\headheight}{14pt}
  \fancyhf{}
  \fancyhead[R]{\small\thepage}
  \renewcommand{\headrulewidth}{0pt}
  \newcommand{\secao}[1]{\clearpage\section{#1}}
  \newcommand{\fontepropria}{\par\smallskip\noindent{\footnotesize\textit{Fonte: elaboração com base no lote experimental final da equipe (2026).}}\par}
  \newcommand{\pendencia}[1]{\par\smallskip\noindent{\small\textbf{A completar antes da entrega:} #1}\par}
  % Campos que dependem da equipe; editar aqui antes da submissão.
  \newcommand{\urlvideo}{\textit{Pendente de gravação e publicação.}}
  \newcommand{\commitfinal}{\textit{A preencher após consolidar os artefatos.}}
  \newcommand{\titulotrabalho}{GENRESEARCH-SERVIDOR: MEMÓRIA, ARMAZENAMENTO, FILAS, CACHE E SINCRONIZAÇÃO}
  \newcommand{\integrantes}{Bruno Amancio Ferreira\\Christian Will Silva Santos Nunes\\Iasmin Victoria Teixeira Barreto\\Marcos Vinícius Dantas Aguiar\\Pedro César Figueiredo Carneiro\\Sibele Oliveira Cruz\\Silas Santos da Silva}
  % Gráficos vetoriais incorporados; valores são médias das três repetições.
  \pgfplotsset{grafico/.style={ybar,width=0.46\textwidth,height=4.15cm,
   ymin=0,bar width=16pt,axis lines=left,enlarge x limits=0.35,
   xtick=data,tick label style={font=\footnotesize},label style={font=\footnotesize},
   title style={font=\small},nodes near coords,
   every node near coord/.append style={font=\scriptsize},
   scaled y ticks=false,grid=major,grid style={gray!18},
   /pgf/number format/use comma}}

  \begin{document}
  \hypersetup{pageanchor=false}
  \begin{titlepage}
  \thispagestyle{empty}
  \begin{center}
  \textbf{UNIVERSIDADE FEDERAL DE SERGIPE}
  \vspace{1.2cm}

  \textbf{ATIVIDADE 2 - SISTEMAS OPERACIONAIS (2026.2)}
  \vfill
  {\bfseries\integrantes}
  \vfill
  {\large\textbf{\titulotrabalho}}
  \vfill
  \textbf{SERGIPE}\\\textbf{2026}
  \end{center}
  \end{titlepage}

  \begin{titlepage}
  \thispagestyle{empty}
  \begin{center}
  \integrantes
  \vfill
  {\large\textbf{\titulotrabalho}}
  \vfill
  \begin{flushright}
  \begin{minipage}{0.55\textwidth}
  Relatório técnico apresentado à Universidade Federal de Sergipe como parte da Atividade 2 da disciplina Sistemas Operacionais, período 2026.2.
  \end{minipage}
  \end{flushright}
  \vfill
  SERGIPE\\2026
  \end{center}
  \end{titlepage}

  \hypersetup{pageanchor=true}
  \pagenumbering{arabic}
  \thispagestyle{empty}
  \begin{center}\textbf{RESUMO}\end{center}
  Este relatório apresenta a implementação e a análise do GenResearch-Servidor, um pipeline inspirado em recuperação aumentada por geração (RAG). O objetivo foi relacionar memória, armazenamento, leitura de arquivos, filas, cache e sincronização ao desempenho e à correção de uma aplicação. O recorte utiliza documentos sintéticos, vetores lexicais simulados, recuperação por similaridade e resposta com espera simulada, sem inferência de modelo de linguagem. Em Ubuntu/WSL~2, foram executadas onze condições com três repetições, totalizando 720 consultas concluídas e 72 falhas intencionais. Compararam-se leitura convencional e mapeada, um a quatro consumidores, três tamanhos de corpus e capacidades de cache. Na carga pequena, o uso de \texttt{mmap} reduziu a janela média de 0,454 para 0,395~s, com memória residente semelhante. A vazão passou de 26,93 para 74,03 consultas por segundo entre um e quatro consumidores. O cache de capacidade três atingiu 31,94\% de acerto e reduziu a janela média em aproximadamente 21,5\%. Uma condição de corrida isolada duplicou o cálculo de uma mesma chave; a coordenação por chave eliminou a duplicação. Não houve swap nas amostras nem evidência suficiente de thrashing. As conclusões são limitadas pelos dados pequenos, pela simulação da geração, pela amostragem e pela integração de arquivos do WSL. O estudo demonstra a relação entre organização dos dados, coordenação concorrente, desempenho e correção, preservando as limitações de generalização.

  \vspace{\baselineskip}
  \noindent\textbf{Palavras-chave:} Sistemas Operacionais. Memória virtual. Mapeamento de arquivos. Concorrência. Cache.

  \clearpage
  \thispagestyle{empty}
  \begin{singlespace}\small\tableofcontents\end{singlespace}

  % Oito páginas de conteúdo, além dos quatro elementos iniciais.
  \secao{Contribuições, introdução e problema}
  \begin{table}[H]
  \centering
  \caption{Distribuição de funções entre os integrantes}
  \label{tab:contribuicoes}
  \begin{singlespace}\small
  \begin{tabularx}{\textwidth}{>{\raggedright\arraybackslash}p{0.49\textwidth}X}
  \toprule\textbf{Integrante} & \textbf{Função atribuída}\\\midrule
  Bruno Amancio Ferreira & Coordenação do trabalho; introdução, objetivos e organização da entrega.\\
  Christian Will Silva Santos Nunes & Arquitetura do GenResearch; pipeline, dados e demonstração da aplicação.\\
  Iasmin Victoria Teixeira Barreto & Ambiente experimental, configurações e metodologia das medições.\\
  Marcos Vinícius Dantas Aguiar & Sincronização do cache; demonstração da corrida, correção e falhas controladas.\\
  Pedro César Figueiredo Carneiro & Processos, threads e fila; análise de read/mmap e chamadas de sistema.\\
  Sibele Oliveira Cruz & Consolidação de métricas, tabelas e gráficos; análise de desempenho e cache.\\
  Silas Santos da Silva & Discussão conceitual, limitações, conclusão, referências e revisão final.\\\bottomrule
  \end{tabularx}
  \end{singlespace}
  \end{table}
  A tabela distribui as responsabilidades técnicas e editoriais entre os sete integrantes, em alinhamento com os temas de participação no vídeo. As funções orientam a preparação, a apresentação e a revisão das respectivas partes do trabalho.

  Na atividade anterior, o grupo observou a inferência local do Gemma~3. Nesta atividade, o recorte muda para um serviço simplificado inspirado em RAG, com controle de arquivos, fila e cache. As medições anteriores de Ollama, ARM64 e Open WebUI não são resultados do GenResearch.

  A pergunta norteadora é: \textit{como memória principal, memória virtual, armazenamento, caches, filas compartilhadas e mecanismos de sincronização afetam o desempenho, a estabilidade e a correção de uma aplicação de IA generativa ou de RAG?}

  O objetivo é implementar um fluxo observável, comparar configurações de acesso e concorrência, demonstrar uma falha controlada e explicar sua correção. O foco está nos custos e mecanismos de Sistemas Operacionais, sem avaliar qualidade de texto de um modelo.

  \noindent\textbf{Repositório:} \url{https://github.com/4rwill/SO_UFS_2026_2_Bruno_Amancio_Christian_Will}.\\
  \textbf{Vídeo:} \urlvideo\\
  \textbf{Commit final:} \commitfinal

  \secao{Ambiente experimental e arquitetura}
  O inventário foi coletado em 10/10/2026 às 15:44:59, fuso $-03{:}00$ (America/Sao\_Paulo). Os valores descrevem a instância Linux, sem representar todo o Windows \cite{projeto}.
  \begin{table}[H]
  \centering\caption{Ambiente registrado}\label{tab:ambiente}
  \begin{singlespace}\small
  \begin{tabularx}{\textwidth}{>{\raggedright\arraybackslash}p{0.23\textwidth}X}
  \toprule\textbf{Item} & \textbf{Observação}\\\midrule
  Sistema & Ubuntu 26.04.1 LTS; WSL~2 sobre Windows.\\
  Kernel & \texttt{6.18.40.1-microsoft-standard-WSL2}.\\
  CPU & Intel Core i5-11260H; 6 núcleos, 12 CPUs lógicas.\\
  RAM/swap & 3,7~GiB RAM; cerca de 3,1~GiB disponíveis no inventário; swap de 1~GiB, 0~B usados.\\
  Arquivos & Projeto em \texttt{/mnt/c}, via 9p/DrvFS; volume de 476~GiB, 323~GiB disponíveis.\\
  Disco Linux & \texttt{/dev/sdd}, ext4; 1007~GiB virtuais apresentados, sem reserva física equivalente garantida.\\
  GPU & NVIDIA GTX~1650; 4096~MiB VRAM; não utilizada.\\
  Aplicação & GenResearch 0.1.0; Python 3.14.4; biblioteca padrão.\\
  Ferramentas & sysstat 12.7.7; strace 6.19; GNU time 1.9; Matplotlib 3.10.8 instalado após as medições.\\\bottomrule
  \end{tabularx}\end{singlespace}
  \end{table}
  \fontepropria
  O programa é executado no terminal. O comando \texttt{serve} processa uma carga local finita, sem servidor HTTP. Seu fluxo é: documentos, fragmentação e vetorização, persistência, produtores, fila limitada, consumidores, cache, recuperação em caso de miss, contexto, resposta simulada e logs.

  \texttt{MemFree} indica páginas livres; \texttt{MemAvailable} estima memória disponibilizável sem swap, considerando páginas recuperáveis. O cache de páginas pode acelerar releituras e não é necessariamente desperdício. RSS mede páginas residentes; VSZ inclui reservas e mapeamentos virtuais. Threads compartilham o espaço de memória: não se deve somar o RSS repetido em suas linhas \cite{proc,threads}.

  O registro aponta o commit base \texttt{e649deb2855910f8aa890bb6858d7fb744e5ce8c}, com alterações locais. Manifestos e hashes identificam os arquivos medidos; esse identificador não substitui o commit final da entrega.

  \secao{Dados, fila e estratégia de sincronização}
  Os textos são sintéticos, com semente 202602. Fragmentos têm até 256 caracteres e sobreposição de 32; vetores têm 64 dimensões. O SHA-256 dos termos define posições e as contagens são normalizadas em L2. Trata-se de representação lexical simulada, com colisões possíveis, sem equivalência a embeddings semânticos.

  \begin{table}[H]
  \centering\caption{Corpus e índice persistidos}\label{tab:corpus}
  \begin{singlespace}\small
  \begin{tabular}{lrrrr}
  \toprule\textbf{Escala} & \textbf{Docs.} & \textbf{Corpus (B)} & \textbf{Índice (B)} & \textbf{Fragmentos}\\\midrule
  Pequena & 12 & 20.797 & 75.314 & 96\\
  Média & 48 & 83.270 & 302.273 & 384\\
  Maior & 192 & 333.022 & 1.206.322 & 1.536\\\bottomrule
  \end{tabular}\end{singlespace}
  \end{table}
  \fontepropria
  \texttt{corpus.bin} concatena documentos; \texttt{index.json} guarda fragmentos e vetores; \texttt{metadata.json} registra origem, parâmetros, offsets e checksums. O JSON duplica textos para simplificar a implementação. A recuperação usa cosseno, top-$k=3$ e contexto de 512 bytes, com UTF-8 preservado. A resposta identifica o contexto e usa \texttt{sleep} de 0,01~s; não há LLM, pesos ou inferência em GPU.

  Dois produtores inserem consultas em uma \texttt{Queue} de capacidade quatro. Consumidores retiram tarefas; o número varia entre um, dois e quatro. A fila coordena espaço e disponibilidade; timeout gera recusa explícita. Após o fim dos produtores, os consumidores drenam a fila e encerram. A correção não depende de \texttt{qsize()} \cite{queue}.

  O cache LRU possui lock para busca, contadores, substituição e chaves em andamento. A coordenação \textit{single-flight} usa um evento por chave: um trabalhador calcula fora do lock e os demais aguardam resultado ou exceção. A chave combina consulta normalizada, versão dos dados e parâmetros. Falhas não são cacheadas. Um lock separado serializa as linhas e a sequência do log JSONL.

  Na corrida isolada, uma barreira força duas threads a observarem a mesma chave ausente. Sem coordenação, há dois cálculos, um duplicado; com single-flight, há um cálculo e nenhum duplicado. Resultados iguais e encerramento foram confirmados. O processo tem watchdog de 5~s. A falha é duplicação de trabalho; não foram demonstrados deadlock ou starvation.

  \secao{Configurações, metodologia e métricas}
  O lote final \texttt{15c7fb31a435409597342b40f0d70665} contém onze condições com três repetições, em processos novos e ordem fixa. Cada repetição submete 24 consultas: nove textos distintos, com três textos recorrentes em 18 posições e seis sufixos específicos. Duas consultas de aquecimento ficam separadas da medição; o cache da aplicação inicia vazio. C6 ausente não tem aquecimento.

  \begin{table}[H]
  \centering\caption{Matriz de comparações}\label{tab:matriz}
  \begin{singlespace}\small
  \begin{tabularx}{\textwidth}{lX}
  \toprule\textbf{Série} & \textbf{Condição e controle}\\\midrule
  C1/C2 & \texttt{read}/\texttt{mmap}; mesmos arquivos e acesso integral.\\
  C3 & 1/2/4 consumidores; 2 produtores e fila 4.\\
  Corpus & 12/48/192 documentos; corpus e índice crescem juntos.\\
  C4 & Cache LRU com capacidades 0/1/3; mesmas consultas.\\
  C5 & Corpus maior; top-$k$/contexto de 3/512 para 32/8192.\\
  C6 & Índice ausente; depois, processos novos com índice disponível.\\\bottomrule
  \end{tabularx}\end{singlespace}
  \end{table}
  A linha de base usa corpus pequeno, read, dois consumidores, cache zero e geração simulada de 0,01~s. C1 é reutilizada nas séries de consumidores, corpus e cache. Esses três fatores atendem à variação mínima. C5 amplia a carga moderadamente, sem produzir pressão intensa de RAM.

  Vazão é o número de conclusões dividido pela janela do serviço, incluindo lançamento e encerramento. Latência vai da chegada ao término, incluindo admissão e fila. Média, mediana e p95 consideram conclusões; p95 é nearest-rank entre 24 consultas, não entre três repetições. Os resultados agregam média e desvio padrão amostral; não foram realizados testes de significância.

  Tempos de ingestão, recuperação, resposta e total são preservados nos artefatos. RSS, VSZ, memória disponível, cache, swap e faltas são amostrados; máximos podem perder picos. Em C1, a média das latências médias foi 122,32~ms; a média das medianas, 139,31~ms. A ocupação média temporal da fila foi 3,28 itens, com máximo quatro. Utilização de trabalhador inclui espera e não é CPU.

  Limites: espaço virtual de 512~MiB, CPU de 30~s e watchdog de 45~s por processo. Não se alteraram swap ou caches do kernel. Manifestos, consultas, recursos, eventos, resumos, horários e hashes foram preservados. A equivalência de conteúdo read/mmap foi verificada nas três repetições.

  \secao{Resultados de leitura e concorrência}
  \subsection{Leitura convencional e mapeada}
  A janela média foi $0{,}454\pm0{,}014$~s em read e $0{,}395\pm0{,}010$~s em mmap, redução observada de 12,9\%. Vazão: $52{,}90\pm1{,}59$ e $60{,}72\pm1{,}57$~req/s; p95: $145{,}31\pm6{,}31$ e $126{,}23\pm5{,}38$~ms. RSS máximo amostrado médio foi semelhante: 34,29 e 34,49~MiB.

  \begin{figure}[H]\centering
  \begin{tikzpicture}\begin{axis}[grafico,symbolic x coords={read,mmap},ylabel={Janela (s)},ymax=0.57]
  \addplot+[fill=blue!45,draw=blue!60!black,error bars/.cd,y dir=both,y explicit] coordinates {(read,0.454) +- (0,0.014) (mmap,0.395) +- (0,0.010)};
  \end{axis}\end{tikzpicture}\hfill
  \begin{tikzpicture}\begin{axis}[grafico,symbolic x coords={1,2,4},ylabel={Vazão (req/s)},xlabel={Consumidores},ymax=95]
  \addplot+[fill=blue!45,draw=blue!60!black,error bars/.cd,y dir=both,y explicit] coordinates {(1,26.93) +- (0,0.22) (2,52.90) +- (0,1.59) (4,74.03) +- (0,2.48)};
  \end{axis}\end{tikzpicture}
  \caption{Janela por modo de leitura e vazão por consumidores. Barras de erro: desvio padrão amostral, três repetições.}\label{fig:leitura-workers}
  \end{figure}
  \fontepropria
  Abertura/mapeamento levou cerca de 0,92/1,32~ms; acesso/cópia/fechamento, 7,24/1,03~ms. O mmap copia integralmente os dados e não é zero-copy \cite{mmap}. Faltas minor médias foram aproximadamente 196.918/98.697; major, 0/24. Não se convertem essas contagens em E/S física ou thrashing sem investigação. Houve aquecimento, sem disco frio verificado; a segunda repetição não foi sempre mais rápida. Não se comparou mmap no corpus maior.

  \subsection{Consumidores e fila}
  Com 1/2/4 consumidores, o p95 foi 257,46/145,31/130,82~ms e a espera média na fila, 129,05/62,17/38,39~ms. O ganho de vazão foi 1,96 vez de um para dois e 1,40 vez de dois para quatro: retorno decrescente, sem queda até quatro ou ponto de saturação identificado.

  Threads sobrepõem leitura e espera; não se demonstra ganho de bytecode limitado por CPU \cite{threads}. O RSS subiu de 31,93 para 50,33~MiB entre um e quatro consumidores. Com quatro, VSZ chegou a 504,67~MiB, próximo do limite virtual de 512~MiB, sem erro de memória observado. Logging, GIL e parsing são custos possíveis, não causas isoladas experimentalmente.

  \secao{Resultados de cache e escala dos dados}
  \subsection{Acertos, substituições e tempo}
  Cache 0/1/3 apresentou hit rate médio de 0/6,94/31,94\% e janelas de 0,454/0,429/0,356~s. A capacidade três reduziu a janela em cerca de 21,5\%; capacidade um, em 5,4\%. As condições ocorreram em momentos distintos, com o mesmo plano de consultas.

  \begin{figure}[H]\centering
  \begin{tikzpicture}\begin{axis}[grafico,symbolic x coords={0,1,3},ylabel={Hit rate (\%)},xlabel={Capacidade do cache},ymax=43]
  \addplot+[fill=blue!45,draw=blue!60!black,error bars/.cd,y dir=both,y explicit] coordinates {(0,0) +- (0,0) (1,6.94) +- (0,2.41) (3,31.94) +- (0,2.41)};
  \end{axis}\end{tikzpicture}\hfill
  \begin{tikzpicture}\begin{axis}[grafico,symbolic x coords={12,48,192},ylabel={RSS (MiB)},xlabel={Documentos},ymax=70]
  \addplot+[fill=blue!45,draw=blue!60!black,error bars/.cd,y dir=both,y explicit] coordinates {(12,34.29) +- (0,1.63) (48,42.48) +- (0,0.32) (192,53.83) +- (0,0.41)};
  \end{axis}\end{tikzpicture}
  \caption{Acertos do cache e memória residente por corpus. RSS: média dos máximos amostrados; barras de erro: desvio padrão, três repetições.}\label{fig:cache-corpus}
  \end{figure}
  \fontepropria
  Capacidade um teve 2/1/2 hits e 19/21/20 substituições; capacidade três, 7/8/8 hits e 12/11/11 substituições. Seguidores single-flight contam como miss/coalesced, sem valor pronto no instante da busca. Miss rate é $1-\text{hit rate}$, sem equivaler a falha de execução. O p95 com cache um foi $152{,}23\pm12{,}63$~ms, acima da linha de base: menor janela não garante menor cauda. A estimativa bruta de tempo evitado não desconta espera/logging e não é economia líquida.

  \subsection{Corpus, contexto e paginação}
  Ao passar de 12 para 48 e 192 documentos, o p95 foi de 145,31 para 169,30 e 317,18~ms; vazão, de 52,90 para 45,44 e 25,25~req/s. Texto e índice crescem juntos, sem isolamento de seus custos. O lote de 24 consultas não foi variado.

  C5 ampliou top-$k$/contexto no corpus maior: RSS foi $52{,}59\pm2{,}41$~MiB e p95, $311{,}95\pm8{,}36$~ms, sem aumento claro frente a 53,83~MiB e 317,18~ms. Carregamento integral e sleep fixo limitam a sensibilidade ao contexto. Swap permaneceu zero nas amostras, sem evidência suficiente de pressão intensa ou thrashing.

  \secao{Correção, chamadas de sistema e discussão}
  O lote totaliza 792 consultas: 720 concluídas e 72 falhas intencionais de índice ausente. Não houve perdas, duplicações de IDs, recusas inesperadas ou threads de trabalho restantes. A duplicação de IDs é diferente da duplicação de cálculo da corrida isolada. C6-recovery concluiu 72 consultas em processos novos com índice disponível, com p95 médio de 147,37~ms; isso não comprova retentativa automática no mesmo processo.

  Na validação funcional de saturação, fora da matriz, três consultas concluíram e 27 foram recusadas explicitamente. A falha da corrida foi corrigida de dois para um cálculo, com resultados iguais. Locks e eventos acrescentam coordenação e espera, mas permitem preservar consistência sem manter lock global durante o cálculo.

  A coleta complementar separada usou sleep de 0,05~s e concluiu 24 consultas por modo. \texttt{ps} e \texttt{/proc} mostraram seis threads: principal, amostrador, dois produtores e dois consumidores. Seus tempos não se misturam à matriz.
  \begin{table}[H]\centering
  \caption{Chamadas selecionadas no rastreamento complementar}\label{tab:syscalls}
  \begin{singlespace}\small
  \begin{tabularx}{\textwidth}{lrrX}
  \toprule\textbf{Chamada} & \textbf{read} & \textbf{mmap} & \textbf{Interpretação}\\\midrule
  \texttt{clone3} & 5 & 5 & Criação de threads.\\
  \texttt{futex} & 6.888 & 6.818 & Espera/sinalização.\\
  \texttt{read} & 720 & 616 & Leitura em descritores.\\
  \texttt{openat} & 511 & 505 & Abertura e tentativas.\\
  \texttt{mmap} & 273 & 273 & Mapeamentos globais.\\
  \texttt{write} & 300 & 300 & Escrita, incluindo logs.\\\bottomrule
  \end{tabularx}\end{singlespace}
  \end{table}
  \fontepropria
  O filtro de corpus/índice encontrou 0/48 mapeamentos nos modos read/mmap. Contagens globais iguais não indicam comportamento igual nos arquivos alvo. Retornos EAGAIN/ETIMEDOUT em futex não são falhas de consulta.

  \texttt{/proc/PID/io} não existia no ambiente: bytes físicos e operações por processo permanecem indisponíveis, não zero. Tamanhos e bytes lógicos não substituem E/S física. \texttt{vmstat}/\texttt{iostat} observam Linux/dispositivos, sem atribuição exclusiva ao SSD Windows.

  Em RAG, documentos, vetores, índices e contextos competem por recursos. Cache de respostas e cache de páginas atuam em níveis distintos. Checkpoints, GPU e execução distribuída adicionariam transferência e sincronização, mas não foram medidos. Auditorias conferiram 33 execuções, 364 hashes de entrada da análise e 51 da coleta complementar; 19 testes passaram na etapa da matriz.

  \secao{Limitações, conclusão e referências}
  As limitações são corpus pequeno sintético, hashing com colisões, resposta com sleep fixo, JSON/cópias, logs síncronos, três repetições em ordem fixa, aquecimento, máximos apenas amostrados, limite virtual e WSL/DrvFS. Faltam E/S física por processo, comparação mmap em corpus maior e reprodução comprovada em clone limpo. Não se avaliaram qualidade textual, inferência real ou GPU.

  Conclui-se que corpus/índice maiores acompanharam maior RSS e latência; mais consumidores aumentaram vazão e reduziram espera; cache maior ampliou reutilização; sincronização por chave evitou trabalho duplicado. O mmap reduziu a janela nesta carga, sem superioridade geral demonstrada. Recomenda-se testar acesso parcial, corpus maior, ordem randomizada e geração real ou custo computacional calibrado.

  \noindent\textbf{Uso de IA.} O projeto registra assistência do Codex em implementação, auditoria e análise; esta redação e adaptação em LaTeX também tiveram assistência de IA. A equipe deve revisar valores, interpretação e o registro das contribuições conforme as funções da Tabela~\ref{tab:contribuicoes}. Antes de submeter, preencher o commit final e a URL do vídeo, também exigida no README e em \texttt{VIDEO.md} ou \texttt{video.txt}.

  \begingroup
  \renewcommand{\section}[2]{}
  \begin{singlespace}\footnotesize
  \begin{thebibliography}{9}
  \bibitem{projeto} EQUIPE GENRESEARCH. \textbf{Código, configurações e resultados}. 2026. \url{https://github.com/4rwill/SO_UFS_2026_2_Bruno_Amancio_Christian_Will}. Fontes locais: documentação; lote \texttt{15c7fb31a435409597342b40f0d70665}; tabelas e evidências. Consulta: 14 out. 2026.
  \bibitem{threads} PYTHON SOFTWARE FOUNDATION. \textbf{threading: Thread-based parallelism}. \url{https://docs.python.org/3.14/library/threading.html}. Acesso: 14 out. 2026.
  \bibitem{queue} PYTHON SOFTWARE FOUNDATION. \textbf{queue: A synchronized queue class}. \url{https://docs.python.org/3.14/library/queue.html}. Acesso: 14 out. 2026.
  \bibitem{proc} THE LINUX KERNEL DOCUMENTATION. \textbf{The /proc Filesystem}. \url{https://www.kernel.org/doc/html/latest/filesystems/proc.html}. Acesso: 14 out. 2026.
  \bibitem{mmap} PYTHON SOFTWARE FOUNDATION. \textbf{mmap: Memory-mapped file support}. \url{https://docs.python.org/3.14/library/mmap.html}. Acesso: 14 out. 2026.
  \bibitem{enunciado} UNIVERSIDADE FEDERAL DE SERGIPE. \textbf{Sistemas Operacionais: Atividade 2 (AV2)}. Enunciado fornecido pela equipe, 2026.
  \end{thebibliography}\end{singlespace}
  \endgroup
  \end{document}
